% C01.09/C01.11 application; concise practice recalls the single theory homes.
\subsection{হিসাব: আগে একক, তারপর সূত্র}
প্রেরিত ডেটা $D$ bit এবং সময় $t$ s হলে ডেটা পাঠানোর পর্যবেক্ষিত গড় হার $R=D/t$ bps। নির্দিষ্ট হারে ডেটা পাঠানোর আদর্শ \textbf{সময় (Transmission Time)} হলো $t=D/R$। অন্যান্য বিলম্ব এবং নিয়ন্ত্রণতথ্য পাঠানোর মতো বাড়তি কাজ (Overhead) থাকলে মোট সময় বাড়বে।
\chTwoPair{\textbf{উদাহরণ ১: গড় হার}\par
5 সেকেন্ডে 9000 bit ডেটা পাঠানো হলে গড় হার $R=9000/5=1800$ bps। আবার 5 সেকেন্ডে 50,000 bit পাঠানোর হিসাবে হার হবে $R=10,000$ bps $=10$ kbps।}{\textbf{উদাহরণ ২: Byte থেকে bit}\par
500 Byte ডেটা দুই সেকেন্ডে পাঠানো হলে $D=500\times8=4000$ bit। সুতরাং গড় হার $R=4000/2=2000$ bps। একই হার বাইটে প্রকাশ করলে 250 Byte/s হয়; 250 bps হয় না।}
\begin{mistakebox}{গড় হার ও সর্বোচ্চ বহনক্ষমতার পার্থক্য}
বাস্তবে পাওয়া $D/t$ একটি পর্যবেক্ষিত গড় হার; তা থেকে সংযোগের সর্বোচ্চ বহনক্ষমতা জানা যায় না। প্রশ্নে সংযোগ সর্বোচ্চ হারে চলেছে বলা থাকলে সেই শর্তে বহনক্ষমতা নির্ণয় করা যায়। হিসাবের আগে মিনিটকে সেকেন্ডে এবং বাইটকে বিটে রূপান্তর করতে হবে। ডেটা হারের এককগুলোর ভিত্তিও মিলিয়ে নিতে হবে।
\end{mistakebox}

\subsection{বাস্তব সমস্যা: স্থানীয় সংযোগ দ্রুত, কিন্তু আপলোড ধীর}
একজন শিক্ষার্থীকে কলেজে 12 MB-এর একটি ভিডিও পাঠাতে হবে। স্থানীয় সংযোগের ক্ষমতা 100 Mbps, ইন্টারনেটে আপলোডের সংযোগ 3 Mbps এবং অপর প্রান্তের গ্রহণক্ষমতা 20 Mbps। পাঠাতে কত সময় লাগবে, তা জানতে শুধু স্থানীয় সংযোগের হার দেখলে চলবে না।
\noindent\begin{minipage}{\linewidth}\centering
\begin{tikzpicture}[node distance=8mm]
\node[chTwoNode,text width=2cm] (a) {Computer};
\node[chTwoNode,text width=2cm,right=of a] (b) {Local Link\\100 Mbps};
\node[chTwoNode,text width=2cm,right=of b,draw=chTwoRose] (c) {Upload Link\\3 Mbps};
\node[chTwoNode,text width=2cm,right=of c] (d) {Receiver\\20 Mbps};
\draw[chTwoArrow] (a)--(b);\draw[chTwoArrow] (b)--(c);\draw[chTwoArrow] (c)--(d);
\end{tikzpicture}\par\raggedright
\captionof{figure}{এই পথে 3 Mbps অংশটি ডেটা হার সীমিত করে; স্থানীয় 100 Mbps পুরো পথের হার নয়}
\end{minipage}
\chTwoPair{\textbf{সমাধান:} এই পথে বটলনেক হলো 3 Mbps আপলোড সংযোগ। দশমিক MB ধরে $D=12\times8=96$ Mbit। আদর্শ প্রেরণসময় $96/3=32$ সেকেন্ড। বাস্তবে থ্রুপুট কম হলে আরও সময় লাগবে।}{\textbf{ব্যবহার:} বড় অ্যাসাইনমেন্ট পাঠানো বা ক্লাউডে ব্যাকআপ রাখার সময় পরিকল্পনা করতে আপলোডের হার জানতে হয়। একই সময়ে বড় ফাইল আপলোড করলে ভিডিও ক্লাসের ডেটা সময়মতো পৌঁছাতে বাধা হতে পারে।}
এমন সমস্যায় অপ্রয়োজনীয় আপলোড ও ডাউনলোড বন্ধ করে সংযোগ পরীক্ষা করা যায়। Wi-Fi সংকেত ও তারের সংযোগও দেখতে হয়। প্রয়োজনে ভিডিওর রেজল্যুশন কমানো যেতে পারে। শুধু 4G/5G পরিচয় বা বড় Mbps মান দেখে সরাসরি অনলাইন ক্লাসের ভালো মান নিশ্চিত করা যায় না।

% Historical context retained from old main file; no invented board frequency claim.
\begin{keypointbox}{সম্পূরক: যোগাযোগের বিকাশ}
\chTwoPair{\textbf{টেলিফোন: Bell ও ১৮৭৬}\par
Alexander Graham Bell ৭ মার্চ ১৮৭৬ সালে টেলিফোন-সংক্রান্ত মার্কিন পেটেন্ট পান। ১০ মার্চ তাঁর পরীক্ষায় বোধগম্য কথা অপর প্রান্তে পৌঁছায়। টেলিফোনের ইতিহাসে তাই Bell-এর নামের সঙ্গে ১৮৭৬ সাল গুরুত্বপূর্ণ। এই সাল পরবর্তী সেলুলার হ্যান্ডসেট আবিষ্কারের সাল নয়।}{\textbf{তথ্যতত্ত্ব: Shannon ও ১৯৪৮}\par
Claude Shannon-এর ১৯৪৮ সালের গবেষণা যোগাযোগে তথ্যের পরিমাণ, নয়েজ এবং চ্যানেলের সীমা বিশ্লেষণের ভিত্তি গড়ে দেয়। তাঁকে তথ্যতত্ত্ব বা Information Theory-এর জনক বলা হয়। টেলিফোনের যন্ত্র, তথ্যতত্ত্ব এবং ইন্টারনেটের প্রোটোকলের বিকাশ পৃথক ঘটনা।}
রেডিওর বিকাশের ইতিহাস বেতার যোগাযোগের অংশে এবং Ethernet/TCP-IP-এর ইতিহাস কম্পিউটার নেটওয়ার্কের অংশে রয়েছে।
\end{keypointbox}
